\chapter{Derivations}\label{app:deriv}

\section{Scaling of the amplitude spectrum}\label{app:scale}
Let $x[n]=A\cos(\theta n)$ with $\theta=2\pi f_0/\fs$, and let the window be $w[n]$, $n=0,\dots,N-1$. Write
$x[n]=\tfrac A2\bigl(e^{\jj\theta n}+e^{-\jj\theta n}\bigr)$. The windowed DFT at a bin with $2\pi k/N=\theta$ is
\begin{align}
  X_w[k] &= \sum_{n}w[n]\,x[n]\,e^{-\jj\theta n}
  = \frac A2\sum_{n}w[n]\;+\;\frac A2\sum_{n}w[n]\,e^{-\jj2\theta n}.
\end{align}
The second sum is the window's transform at twice the tone frequency, far outside the main lobe, and is negligible for a
tapered window. Hence $\abs{X_w[k]}\approx\tfrac A2\sum_nw[n]$, which is Eq.~\eqref{eq:ampscale}. If only a subset
$\mathcal V$ of the samples is valid and the others are set to zero, the first sum runs over $\mathcal V$ and the scale becomes
$\sum_{n\in\mathcal V}w[n]$, the rule used by the application.

\section{Hann window constants}\label{app:hann}
For $w[n]=\tfrac12\bigl(1-\cos\frac{2\pi n}{N-1}\bigr)$ and large $N$, the mean over a full period of the cosine is zero, so
$\sum w\approx N/2$, giving a coherent gain of $\tfrac12$. For the energy, with $c=\cos\frac{2\pi n}{N-1}$ averaging to 0 and $c^2$ to $\tfrac12$,
\[
  \sum w^2\approx N\,\overline{\tfrac14\bigl(1-2c+c^{2}\bigr)}=\tfrac N4\bigl(1+\tfrac12\bigr)=\tfrac{3N}{8}.
\]
The equivalent noise bandwidth is therefore
\[
  \mathrm{ENBW}=\frac{N\sum w^2}{(\sum w)^2}=\frac{N\cdot 3N/8}{N^2/4}=1.5\ \text{bins},
\]
matching Table~\ref{tab:windows}.

\section{Zero-phase filtering}\label{app:zp}
Let a filter have impulse response $h$ and frequency response $H(\omega)$. Filtering forward gives $Y_1=H X$. Time reversal of a real
sequence turns its transform into the complex conjugate (up to a linear phase that is undone by the second reversal), so the
backward pass multiplies by $H^{*}$ and the final reversal restores the time direction: $Y=H^{*}H\,X=\abs{H}^{2}X$. The
magnitude response in decibels is $20\log_{10}\abs{H}^{2}=2\cdot20\log_{10}\abs{H}$, twice that of one pass, so the $-3$~dB corner of one
pass is a $-6$~dB corner of the zero-phase filter.

\section{The recursive averages}\label{app:rec}
Rearranging Eq.~\eqref{eq:rec}: $\mathrm{STA}_i=(1-\tfrac1N)\mathrm{STA}_{i-1}+\tfrac1N c[i]$. Unrolling,
\begin{equation}
  \mathrm{STA}_i = \frac1N\sum_{k=0}^{i}\Bigl(1-\frac1N\Bigr)^{k}c[i-k]+\Bigl(1-\frac1N\Bigr)^{i+1}\mathrm{STA}_{-1}.
\end{equation}
The weights sum to $1-(1-1/N)^{i+1}\to1$ and, for large $N$, $(1-1/N)^{k}\approx e^{-k/N}$: an exponential window of time
constant $N$ samples. The initial-condition term decays with the same time constant, which is why a warm-up of $N_l$ samples is
needed.

\section{Detection delay}\label{app:delay}
For a step of the CF power from $P_0$ to $P_1$ at $t=0$ the continuous-time analogue of~\eqref{eq:rec} is
$\dot{S}=(P_1-S)/\tau_s$ with $S(0)=P_0$, whose solution is $S(t)=P_1+(P_0-P_1)e^{-t/\tau_s}$. Setting
$S(t_d)=r_{\mathrm{on}}P_0$ and solving for $t_d$,
\[
  e^{-t_d/\tau_s}=\frac{P_1-r_{\mathrm{on}}P_0}{P_1-P_0}
  \;\Rightarrow\;
  t_d=-\tau_s\ln\frac{P_1-r_{\mathrm{on}}P_0}{P_1-P_0},
\]
which is Eq.~\eqref{eq:delay}. The assumption that the LTA is constant during $t_d$ holds when $t_d\ll\tau_l$; for $\tau_l=30$~s and
$\tau_s=1$~s it is good for all events strong enough to trigger within a few seconds.

\section{The $F$ distribution of the ratio}\label{app:f}
If the CF is the square of zero-mean Gaussian noise of variance $\sigma^{2}$ filtered to bandwidth $B$, the average of $c$ over $T$ seconds is
$\sigma^{2}\chi^{2}_{\nu}/\nu$ with $\nu\approx2BT$ degrees of freedom (the number of independent samples is about $2BT$, by the
sampling theorem). The STA and LTA are independent only if their windows do not overlap, which the classic form approximately
satisfies, and then $R=\dfrac{\chi^{2}_{\nu_s}/\nu_s}{\chi^{2}_{\nu_l}/\nu_l}\sim F_{\nu_s,\nu_l}$. The tail probabilities of
Table~\ref{tab:null} are $1-F_{\nu_s,\nu_l}(r)$ with $\nu_s=18$, $\nu_l=540$. They apply to \emph{one} independent evaluation;
the number of false alarms per day also depends on how many independent evaluations occur and on how exceedances cluster,
which is why the table reports simulated counts alongside the tail.

\section{Error of min/max decimation}\label{app:mm}
Let a bin contain samples $x[n]$, $n\in b$, and let $\hat x$ denote the drawn polyline through the kept extremes. Every sample of the bin lies
between its minimum and maximum, so the vertical distance between any sample and the plotted bin is at most
$x^{\max}_b-x^{\min}_b$, and the horizontal position of the extremes is exact. Therefore the plotted envelope
\[
  \bigl[\,x^{\min}_b,\ x^{\max}_b\,\bigr]
\]
equals the true envelope of the bin, and the only information lost is the order of the samples between the extremes.
